\documentclass[11pt]{ctexart}
\usepackage[a4paper,margin=1in]{geometry}
\usepackage{graphicx}
\usepackage{xcolor}
\usepackage{hyperref}
\definecolor{refgray}{gray}{0.45}
\title{\textbf{市场最新观点汇总}\\[0.4em]{\large 能源冲击退潮与结构性分化并行，黄金、日股、中国消费各寻新定价锚}}
\author{KC桌面 自动生成}
\date{260621}
\begin{document}
\maketitle
\tableofcontents
\newpage
\section{一页摘要}
\begin{itemize}
  \item 亚洲通胀最猛烈阶段已过，但核心通胀路径、工资传导和政策退出节奏的分化将成为未来6-12个月核心变量。
  \item 黄金结构性支撑仍在，但短期上行需ETF资金流重新激活，央行购金和地缘缓和构成底部支撑。
  \item 中东局势缓和利好日本个人护理和化妆品板块成本端，但造纸板块面临提价难度加大的矛盾。
  \item 日本制药板块正经历资金流向的结构性出清，下半年催化剂事件或成反转关键。
  \item 中国美妆市场618大促显示高端化是结构性力量，外资高端品牌和部分国货高端品牌领跑。
\end{itemize}
\section{宏观与利率：亚洲通胀见顶，结构性分化开启}
\textbf{亚洲通胀的能源冲击阶段已经过去，但各经济体在核心通胀粘性、工资传导和政策空间上的分化将重新定义资产定价。}

\begin{itemize}
  \item GS将布伦特原油Q4预测从90美元/桶下调至80美元/桶，主要基于霍尔木兹海峡恢复通航前景改善，油价已从4月高点140美元/桶回落至6月约80美元/桶。
  \item 亚洲整体CPI中位数已升至约3\%，但核心通胀总体温和，多数经济体CPI仍略高于目标区间。
  \item 菲律宾、泰国等低收入经济体过去三个月CPI年化超10\%，但5月环比已明显放缓；高收入经济体（日本除外）工资通胀偏软。
  \item 能源价格回落尚未完全传导至终端消费品，进口价格滞后效应意味着下月通胀数据预计明显改善。
  \item 中国CPI通胀仍处于低位，PPI和进口价格指数显示上游压力缓解，但零售能源价格调整存在滞后。
  \item 韩国、日本等经济体进口价格同比在3-4月剧烈波动，部分5月仍为负值，反映能源成本传导的不对称性。
  \item 印度、印尼等经济体通胀压力相对可控，但零售能源价格仍高于战前水平，补贴政策缓冲了部分冲击。
  \item 通胀预测风险已转向下行，但中东局势仍是最大变量，供给侧尾部风险收窄但需求侧韧性未被充分定价。
\end{itemize}
\begin{figure}[htbp]
\centering
\includegraphics[width=0.88\linewidth]{figures/fig_003.jpg}
\caption{Exhibit 3 - Exhibit 3: Oil prices are down sharply over the past month}
\end{figure}
\begin{figure}[htbp]
\centering
\includegraphics[width=0.88\linewidth]{figures/fig_005.jpg}
\caption{Exhibit 5 - Exhibit 5: Our base case oil forecast now has Brent at \textbackslash{}\$80/bbl in Q4 Monthly Brent oil price}
\end{figure}
\begin{figure}[htbp]
\centering
\includegraphics[width=0.88\linewidth]{figures/fig_006.jpg}
\caption{Exhibit 6 - Exhibit 6: Median headline CPI inflation up to 3\% regionally}
\end{figure}
\begin{figure}[htbp]
\centering
\includegraphics[width=0.88\linewidth]{figures/fig_007.jpg}
\caption{Exhibit 7 - Exhibit 7: Core inflation generally has been well-behaved so far}
\end{figure}
\begin{figure}[htbp]
\centering
\includegraphics[width=0.88\linewidth]{figures/fig_009.jpg}
\caption{Exhibit 11 - Exhibit 11: China CPI inflation}
\end{figure}
\begin{figure}[htbp]
\centering
\includegraphics[width=0.88\linewidth]{figures/fig_014.jpg}
\caption{Exhibit 16 - Exhibit 16: China retail energy prices China weighted average fuel price}
\end{figure}
\begin{figure}[htbp]
\centering
\includegraphics[width=0.88\linewidth]{figures/fig_017.jpg}
\caption{Exhibit 19 - Exhibit 19: Korea CPI contribution breakdown}
\end{figure}
\begin{figure}[htbp]
\centering
\includegraphics[width=0.88\linewidth]{figures/fig_029.jpg}
\caption{Exhibit 31 - Exhibit 31: India CPI contribution breakdown}
\end{figure}
{\small\color{refgray}\textbf{References:} [R001] GS：亚洲通胀最猛烈的阶段已经过去，但结构性分化才刚刚开始}

\section{黄金：央行买盘托底，ETF回归是关键变量}
\textbf{黄金结构性支撑坚实，但短期内从当前水平走向5200美元/盎司的关键在于ETF资金流能否重新激活，而非央行购金。}

\begin{itemize}
  \item MS维持黄金看多倾向，但指出5200美元/盎司目标实现难度加大，需要ETF需求配合。
  \item 中东冲突缓和反而利好黄金，因为供给冲击下黄金避险功能受限，缓和期间黄金平均日表现为正回报。
  \item 美联储鹰派立场推高加息预期，提高了持有黄金的机会成本，尤其影响ETF流入。
  \item 历史数据显示，加息后一个月金价平均仍上涨0.84\%，表现并不一致。
  \item 中国央行购金加速，3-5月买入23吨，此前12个月合计仅19吨，官方需求是结构性支撑。
  \item 印度5月上调黄金进口关税以抑制进口，进口量随金价回调而温和。
  \item ETF对利率路径、实际收益率和美元走势高度敏感，目前仍在净卖出状态。
  \item 油价走低若能传导至利率预期，可能改变美联储立场，为黄金提供额外支撑。
\end{itemize}
\begin{figure}[htbp]
\centering
\includegraphics[width=0.88\linewidth]{figures/fig_085.jpg}
\caption{Exhibit 1 - Exhibit 1: Gold has reconnected with real yields}
\end{figure}
\begin{figure}[htbp]
\centering
\includegraphics[width=0.88\linewidth]{figures/fig_087.jpg}
\caption{Exhibit 3 - Exhibit 3: Gold is less effective as a safe haven in supply shocks than demand shocks Annualised Performance Post Crisis (CAGR)}
\end{figure}
\begin{figure}[htbp]
\centering
\includegraphics[width=0.88\linewidth]{figures/fig_088.jpg}
\caption{Exhibit 5 - Exhibit 5: The market pricing for Fed rate hikes has increased dramatically}
\end{figure}
\begin{figure}[htbp]
\centering
\includegraphics[width=0.88\linewidth]{figures/fig_089.jpg}
\caption{Exhibit 7 - Exhibit 7: ETFs have been selling some gold with the Fed on hold}
\end{figure}
\begin{figure}[htbp]
\centering
\includegraphics[width=0.88\linewidth]{figures/fig_090.jpg}
\caption{Exhibit 4 - Exhibit 4: And gold has been performing better on de-escalation days Average Daily Performance Since Ceasefire On Market Events}
\end{figure}
\begin{figure}[htbp]
\centering
\includegraphics[width=0.88\linewidth]{figures/fig_091.jpg}
\caption{Exhibit 6 - Exhibit 6: Real yields are still well above February levels}
\end{figure}
\begin{figure}[htbp]
\centering
\includegraphics[width=0.88\linewidth]{figures/fig_095.jpg}
\caption{Exhibit 11 - Exhibit 11: China's gold purchases have accelerated with the pullback China Central Bank Monthly Gold Purchases (tonnes)}
\end{figure}
\begin{figure}[htbp]
\centering
\includegraphics[width=0.88\linewidth]{figures/fig_096.jpg}
\caption{Exhibit 12 - Exhibit 12: India's imports moderated, with the gold import duty hiked during May to discourage purchases India Gold Imports versus Gold Price}
\end{figure}
{\small\color{refgray}\textbf{References:} [R002] 摩根斯坦利：黄金的下一段行情，关键不在央行，而在ETF能否重新入场}

\section{能源与大宗：油价预期下调，成本传导时滞创造机会}
\textbf{能源价格拐点已至，但成本改善传导至企业利润需要4-6个月，不同行业受益程度和时滞差异显著。}

\begin{itemize}
  \item GS将布伦特原油Q4预测从90美元/桶下调至80美元/桶，霍尔木兹海峡恢复通航前景改善是主因。
  \item 原油价格从4月高点140美元/桶跌至6月约80美元/桶，降幅明显，但炼油产品价格仍高于战前水平。
  \item 亚洲经济体进口价格年初大幅上涨，部分国家通过补贴缓冲零售燃料价格，但企业端成本压力未完全消除。
  \item 韩国、中国、日本等国的原油进口来源已在调整，以应对中东供应风险。
  \item 油价每变动1美元/桶，对Unicharm和花王营业利润影响约5-6亿日元，对狮王影响约1-2亿日元，传导时滞4-6个月。
  \item 油价下跌对日本造纸企业（王子制纸、日本制纸）成本端有利，但可能削弱其7-8月提价计划的议价能力。
  \item 中东局势缓和降低了化妆品包装材料和部分原料的供应短缺风险，但成本占比低，影响有限。
  \item 能源价格回落预计在下月通胀数据中明显体现，但中东局势仍是最大上行风险。
\end{itemize}
\begin{figure}[htbp]
\centering
\includegraphics[width=0.88\linewidth]{figures/fig_003.jpg}
\caption{Exhibit 3 - Exhibit 3: Oil prices are down sharply over the past month}
\end{figure}
\begin{figure}[htbp]
\centering
\includegraphics[width=0.88\linewidth]{figures/fig_004.jpg}
\caption{Exhibit 4 - Exhibit 4: Refined product prices are also down, though they remain notably above pre-war levels}
\end{figure}
\begin{figure}[htbp]
\centering
\includegraphics[width=0.88\linewidth]{figures/fig_005.jpg}
\caption{Exhibit 5 - Exhibit 5: Our base case oil forecast now has Brent at \textbackslash{}\$80/bbl in Q4 Monthly Brent oil price}
\end{figure}
\begin{figure}[htbp]
\centering
\includegraphics[width=0.88\linewidth]{figures/fig_001.jpg}
\caption{Exhibit 1 - Exhibit 1: Energy flows through the Strait of Hormuz have yet to revive... Estimated oil exports through Strait of Hormuz, based on reported vessel count}
\end{figure}
\begin{figure}[htbp]
\centering
\includegraphics[width=0.88\linewidth]{figures/fig_002.jpg}
\caption{Exhibit 2 - Exhibit 2: ...though some Asian economies have secured imports from other sources}
\end{figure}
{\small\color{refgray}\textbf{References:} [R001] GS：亚洲通胀最猛烈的阶段已经过去，但结构性分化才刚刚开始; [R003] JPM：中东局势缓和，日本日化与造纸股迎来结构性成本重估}

\section{日本股市：板块分化加剧，制药出清与消费成本重估}
\textbf{日本股市整体创新高，但板块分化极端：制药板块经历资金结构性出清，而个人护理和化妆品受益于成本改善预期。}

\begin{itemize}
  \item 日经225指数在6月19日创历史新高，但涨幅高度集中于AI和半导体相关板块，非铁金属、电子设备等板块三个月涨幅28-44\%。
  \item 日本制药板块过去一个月下跌5.6\%，三个月下跌11.1\%，而TOPIX同期上涨5.0\%和12.1\%，资金流向挤压是主因。
  \item JPM认为制药板块疲弱并非基本面恶化，而是市场用AI时代估值逻辑重新定价，缺乏明确催化剂的公司将继续承压。
  \item 下半年催化剂事件（10月学术会议、4-6月财报季）可能成为板块反转关键，Kyowa Kirin的KHK4951 2期数据值得关注。
  \item Sumitomo Pharma上周涨6.9\%领跑板块，EHA大会数据提供短期支撑；Chugai提交sparsentan上市申请。
  \item 中东局势缓和利好日本个人护理板块，Unicharm、花王、狮王成本端改善预计在FY2027体现，但需警惕FY2026下半年高价原料库存压力。
  \item 化妆品板块受益于消费者信心回升和游客增加，销售环境改善比成本下降更重要。
  \item 造纸板块面临油价下跌与提价计划之间的矛盾，成本转嫁机制建立可能延迟。
\end{itemize}
\begin{figure}[htbp]
\centering
\includegraphics[width=0.88\linewidth]{figures/fig_097.jpg}
\caption{Figure 12 - Figure 12: New Cases of COVID-19 in Japan (Fixed Point Observation)}
\end{figure}
\begin{figure}[htbp]
\centering
\includegraphics[width=0.88\linewidth]{figures/fig_098.jpg}
\caption{Figure 13 - Figure 13: Cases of Influenza in Japan (Fixed Point Observation)}
\end{figure}
\begin{figure}[htbp]
\centering
\includegraphics[width=0.88\linewidth]{figures/fig_099.jpg}
\caption{Figure 14 - Figure 14: List of TRx Figure 15: Takeda Pharmaceutical (4502) Entyvio IV Number of US Prescriptions (TRx)}
\end{figure}
\begin{figure}[htbp]
\centering
\includegraphics[width=0.88\linewidth]{figures/fig_100.jpg}
\caption{Figure 16 - Figure 16: Takeda Pharmaceutical (4502) Entyvio Pen Number of US Prescriptions (TRx)}
\end{figure}
\begin{figure}[htbp]
\centering
\includegraphics[width=0.88\linewidth]{figures/fig_101.jpg}
\caption{Figure 17 - Figure 17: Takeda Pharmaceutical (4502) TAKHZYRO Number of US Prescriptions (TRx)}
\end{figure}
\begin{figure}[htbp]
\centering
\includegraphics[width=0.88\linewidth]{figures/fig_102.jpg}
\caption{Figure 19 - Figure 19: Astellas Pharma (4503) XTANDI and Competitors' Number of US Prescriptions (TRx)}
\end{figure}
\begin{figure}[htbp]
\centering
\includegraphics[width=0.88\linewidth]{figures/fig_103.jpg}
\caption{Figure 21 - Figure 21: Astellas Pharma (4503) Competitor of VEOZAH Number of US Prescriptions (TRx)}
\end{figure}
\begin{figure}[htbp]
\centering
\includegraphics[width=0.88\linewidth]{figures/fig_104.jpg}
\caption{Figure 18 - Figure 18: Takeda Pharmaceutical (4502) VYVANSE Number of US Prescriptions (TRx)}
\end{figure}
{\small\color{refgray}\textbf{References:} [R003] JPM：中东局势缓和，日本日化与造纸股迎来结构性成本重估; [R004] JPM：日本制药板块正在经历的不是回调，而是结构性出清}

\section{中国消费：618大促揭示高端化结构性力量}
\textbf{中国美妆市场真正的分水岭不是国货与外资，而是高端与大众，高端化正在从趋势变为结构性力量。}

\begin{itemize}
  \item 抖音美妆Top 500品牌整体GMV同比增长约35\%，但头部50个品牌增速约40\%，100名之后品牌增速回落至20\%左右，市场加速集中。
  \item 外资品牌整体完成率约140\%，国产品牌约110\%，但驱动外资的是雅诗兰黛集团（+56\%）和欧莱雅集团高端线。
  \item 雅诗兰黛主品牌增长65\%，海蓝之谜增长54\%；欧莱雅集团中修丽可增长89\%，YSL增长74\%，而巴黎欧莱雅仅增7\%，美宝莲下滑7\%。
  \item 在高端价格带，国货品牌毛戈平表现突出，6月1-18日增速91\%，全周期完成率150\%，类似双十一的控节奏打法。
  \item 林清轩抖音全周期增速188\%，完成率356\%（基数较低）；巨子生物全周期增速15\%，6月加速到44\%。
  \item 珀莱雅全周期增长24\%，但6月下滑5\%；华熙生物和上美集团承压，完成率仅77\%和63\%。
  \item 高端品牌在抖音的增速从5月的20-30\%提升至6月的40-50\%+，中高端/大众线反而从40\%+放缓至20\%左右。
  \item 外资巨头在抖音的增长并非全线复苏，而是集团战略性地将资源向高利润高端线集中。
\end{itemize}
\begin{figure}[htbp]
\centering
\includegraphics[width=0.88\linewidth]{figures/fig_131.jpg}
\caption{Exhibit 2 - Exhibit 2: Tmall+Douyin+JD: Growth of MNC brands has outperformed local brands since 4Q25 Cosmetics Monthly Online GMV YoY by Brand Positioning}
\end{figure}
\begin{figure}[htbp]
\centering
\includegraphics[width=0.88\linewidth]{figures/fig_132.jpg}
\caption{Exhibit 3 - Exhibit 3: Tmall+Douyin+JD: Growth of Domestic Premium brands (i.e. Mao Geping) has outperformed MNC peers Cosmetics Monthly Online GMV YoY by Brand Positioning}
\end{figure}
\begin{figure}[htbp]
\centering
\includegraphics[width=0.88\linewidth]{figures/fig_133.jpg}
\caption{Exhibit 4 - Exhibit 4: Tmall+Douyin+JD: Growth of MNC Mid-to-high end brands has outperformed local peers in 2026 Cosmetics Monthly Online GMV YoY by Brand Positioning}
\end{figure}
\begin{figure}[htbp]
\centering
\includegraphics[width=0.88\linewidth]{figures/fig_134.jpg}
\caption{Exhibit 5 - Exhibit 5: Tmall+Douyin+JD: In mass brands, local has outperformed MNC growth Cosmetics Monthly Online GMV YoY by Brand Positioning Cosmetics Monthly Online GMV YoY by Brand Positioning}
\end{figure}
{\small\color{refgray}\textbf{References:} [R005] GS：618大促揭示中国美妆市场真正的分水岭，不是国货与外资，而是高端与大众}

\section{风险偏好与资金流向：AI主导的集中化与板块轮动}
\textbf{全球市场风险偏好分化，资金高度集中于AI和半导体主题，其他板块面临被动失血，但地缘缓和和成本改善可能触发轮动。}

\begin{itemize}
  \item 日本股市上涨高度集中于AI和半导体，非铁金属、电子设备、玻璃与陶瓷产品三个月涨幅28-44\%，制药板块同期下跌11\%。
  \item 黄金ETF资金流对利率路径高度敏感，美联储鹰派立场压制ETF流入，但央行购金提供结构性支撑。
  \item 中东局势缓和降低了供给冲击风险，油价回落可能改善通胀预期，为风险资产提供支撑。
  \item 亚洲通胀风险转向下行，各国央行政策空间可能逐步打开，但节奏分化。
  \item 中国美妆市场高端化趋势显示消费升级仍在继续，但大众市场增速放缓，品牌集中度提升。
  \item 日本个人护理和化妆品板块受益于成本改善和销售环境好转，可能成为资金轮动方向。
  \item 日本制药板块下半年催化剂事件可能吸引资金回流，但需要明确的数据和事件驱动。
  \item 地缘政治风险（中东、霍尔木兹海峡）仍是最大不确定性，但缓和趋势若持续，将利好风险偏好。
\end{itemize}
\begin{figure}[htbp]
\centering
\includegraphics[width=0.88\linewidth]{figures/fig_003.jpg}
\caption{Exhibit 3 - Exhibit 3: Oil prices are down sharply over the past month}
\end{figure}
\begin{figure}[htbp]
\centering
\includegraphics[width=0.88\linewidth]{figures/fig_005.jpg}
\caption{Exhibit 5 - Exhibit 5: Our base case oil forecast now has Brent at \textbackslash{}\$80/bbl in Q4 Monthly Brent oil price}
\end{figure}
\begin{figure}[htbp]
\centering
\includegraphics[width=0.88\linewidth]{figures/fig_085.jpg}
\caption{Exhibit 1 - Exhibit 1: Gold has reconnected with real yields}
\end{figure}
\begin{figure}[htbp]
\centering
\includegraphics[width=0.88\linewidth]{figures/fig_088.jpg}
\caption{Exhibit 5 - Exhibit 5: The market pricing for Fed rate hikes has increased dramatically}
\end{figure}
\begin{figure}[htbp]
\centering
\includegraphics[width=0.88\linewidth]{figures/fig_089.jpg}
\caption{Exhibit 7 - Exhibit 7: ETFs have been selling some gold with the Fed on hold}
\end{figure}
\begin{figure}[htbp]
\centering
\includegraphics[width=0.88\linewidth]{figures/fig_095.jpg}
\caption{Exhibit 11 - Exhibit 11: China's gold purchases have accelerated with the pullback China Central Bank Monthly Gold Purchases (tonnes)}
\end{figure}
{\small\color{refgray}\textbf{References:} [R001] GS：亚洲通胀最猛烈的阶段已经过去，但结构性分化才刚刚开始; [R002] 摩根斯坦利：黄金的下一段行情，关键不在央行，而在ETF能否重新入场; [R003] JPM：中东局势缓和，日本日化与造纸股迎来结构性成本重估; [R004] JPM：日本制药板块正在经历的不是回调，而是结构性出清; [R005] GS：618大促揭示中国美妆市场真正的分水岭，不是国货与外资，而是高端与大众}

\section{结语}
综合来看，市场正处于能源冲击退潮后的结构性分化阶段：亚洲通胀见顶但分化开启，黄金需ETF回归才能突破，日本股市板块极端分化，中国消费高端化趋势强化。地缘缓和与成本改善为部分板块提供支撑，但资金高度集中于AI主题，其他板块需等待催化剂触发轮动。
\newpage
\section{报告覆盖清单}
\begin{itemize}
  \item [R001] 宏观与利率：亚洲通胀见顶，结构性分化开启 / 能源与大宗：油价预期下调，成本传导时滞创造机会 / 风险偏好与资金流向：AI主导的集中化与板块轮动 - GS：亚洲通胀最猛烈的阶段已经过去，但结构性分化才刚刚开始
  \item [R002] 黄金：央行买盘托底，ETF回归是关键变量 / 风险偏好与资金流向：AI主导的集中化与板块轮动 - 摩根斯坦利：黄金的下一段行情，关键不在央行，而在ETF能否重新入场
  \item [R003] 能源与大宗：油价预期下调，成本传导时滞创造机会 / 日本股市：板块分化加剧，制药出清与消费成本重估 / 风险偏好与资金流向：AI主导的集中化与板块轮动 - JPM：中东局势缓和，日本日化与造纸股迎来结构性成本重估
  \item [R004] 日本股市：板块分化加剧，制药出清与消费成本重估 / 风险偏好与资金流向：AI主导的集中化与板块轮动 - JPM：日本制药板块正在经历的不是回调，而是结构性出清
  \item [R005] 中国消费：618大促揭示高端化结构性力量 / 风险偏好与资金流向：AI主导的集中化与板块轮动 - GS：618大促揭示中国美妆市场真正的分水岭，不是国货与外资，而是高端与大众
\end{itemize}
\section{逐篇报告摘录}
\subsection{[R001] GS：亚洲通胀最猛烈的阶段已经过去，但结构性分化才刚刚开始}
{\small\color{refgray}归类：宏观与利率：亚洲通胀见顶，结构性分化开启 / 能源与大宗：油价预期下调，成本传导时滞创造机会 / 风险偏好与资金流向：AI主导的集中化与板块轮动}

当能源价格从峰值回落、供应链瓶颈逐步疏通，市场对亚洲通胀的焦虑正在快速降温。但GS最新发布的《亚太通胀监测》报告揭示了一个更微妙的图景：通胀压力整体见顶，并不意味着所有经济体都将同步回归平静。真正值得关注的，不是通胀是否回落，而是回落之后，各经济体的价格粘性、政策空间和结构韧性将如何重新定义资产定价和投资逻辑。

这份发布于2026年6月的报告，基于截至5月的CPI/PPI数据以及6月中旬的能源价格走势，给出了一个清晰的主判断：亚洲通胀的能源冲击阶段已经过去，但核心通胀的路径分化、工资传导的不对称性，以及政策退出的节奏差异，将成为未来6至12个月市场定价的核心变量。GS同时下调了布伦特原油Q4预测至80美元/桶，这一调整本身就在释放一个信号——供给侧的尾部风险正在收窄，但需求侧的韧性尚未被充分定价。

1. 能源价格冲击正在退潮，但输入型通胀的滞后效应仍未被完全消化

GS大宗商品团队将布伦特原油Q4预测从90美元/桶下调至80美元/桶，直接原因是霍尔木兹海峡恢复通航的前景改善。报告中的图表清晰显示，原油价格已从4月超过140美元/桶的峰值大幅回落，6月布伦特现货价格已降至约80美元/桶。这是一个重要的拐点信号。

但真正需要推敲的是传导链条的时滞。亚洲经济体在2026年初经历了剧烈的进口价格和生产者价格飙升，尽管部分国家通过隐性或显性补贴缓冲了零售燃料价格的上行，但企业端的成本压力并未完全消除。GS的数据显示，韩国、中国、日本及其他亚洲经济体的进口价格同比变化在3月至4月间出现剧烈波动，部分经济体在5月仍处于负值区间。这意味着，能源价格的回落尚未完全反映到终端消费品价格中。

对投资者而言，这一滞后效应意味着两个层面的含义。其一，短期内核心通胀可能仍有一定惯性，尤其是那些补贴力度较小、价格传导更直接的低收入经济体。其二，一旦滞后效应释放完毕，通胀下行的速度可能快于市场预期。GS在报告中明确表示，通胀预测的风险已转向下行，这正是基于能源价格持续低于峰值的假设。

2. 低收入经济体的通胀压力最为尖锐，但5月数据已出现缓和信号

报告中最值得关注的细节之一，是不同收入水平经济体之间的通胀分化。GS指出，菲律宾和泰国等补贴能力有限的经济体，在过去三个月经季节性调整后的CPI年化通胀率仍超过10\%。这是非常高的水平，反映出这些经济体在面对能源供给冲击时的脆弱性。

但5月的环比数据已经释放出缓和信号。GS特别提到，这些经济体的5月环比通胀率已明显下降。这意味着，最糟糕的阶段可能已经过去。然而，问题在于这些经济体的核心通胀是否也会同步回落。低收入经济体通常食品和能源在CPI篮子中的权重更高，能源价格回落对整体通胀的拉低效果会更直接，但工资传导和本币贬值的压力可能使核心通胀更具粘性。

这一分化对资产配置的意义在于，投资者需要区分“周期性通胀回落”和“结构性通胀中枢上移”。对于菲律宾、泰国等经济体，通胀回落更多是周期性的，但若工资增长或本币贬值形成新的上行压力，央行政策退出的节奏就会更加谨慎。这直接影响到这些市场的利率敏感型资产和汇率预期。

GS在报告中给出了一个容易被忽视的观察：亚太地区的工资通胀整体表现分化，但多数经济体的工资增长稳定甚至偏弱。一个突出的例外是日本。报告明确指出，日本是较高收入经济体中工资偏软的例外——实际上，日本工资通胀表现相对较强。

这一判断与市场对日本央行退出超宽松政策的预期高度相关。如果日本的工资增长持续高于其他发达经济体，日本央行将更有信心实现通胀目标的可持续性，从而加速政策正常化。但GS的报告并未就此展开详细论证，这恰恰是一个值得深入追问的伏笔：日本的工资-通胀螺旋是否已经形成？如果形成，对日本国债收益率曲线和日元汇率的影响将如何传导至亚洲其他市场？

从投资框架的角度看，日本的工资数

[... middle omitted ...]

报告并未完全回答一个关键问题：核心通胀的“良好表现”是否可持续？

这里存在两个尚未被充分讨论的风险。第一，服务价格的粘性。能源价格回落主要影响商品价格，但服务业通胀通常更受工资和租金驱动，调整速度更慢。如果工资增长开始加速，服务价格可能接力能源成为新的通胀推手。第二，供应链重构带来的结构性成本上升。亚洲经济体在过去两年经历了显著的供应链调整，从“效率优先”转向“安全优先”，这种转变往往伴随着更高的生产成本，最终可能以更高的核心通胀中枢形式体现。

GS的报告并未否认这些风险，但也没有给出明确的量化判断。这意味着，投资者在参考这份报告时，需要自行构建对核心通胀路径的情景分析。如果核心通胀在2026年下半年仍维持在3\%以上，亚洲央行的政策转向可能比市场预期的更慢。

5. 报告尚未完全回答的关键问题：通胀回落的斜率将如何影响央行政策路径

GS的报告提供了丰富的通胀数据和趋势判断，但有一个关键问题并未充分展开：通胀回落的斜率是否足以让亚洲央行在2026年下半年转向宽松？

\subsection{[R002] 摩根斯坦利：黄金的下一段行情，关键不在央行，而在ETF能否重新入场}
{\small\color{refgray}归类：黄金：央行买盘托底，ETF回归是关键变量 / 风险偏好与资金流向：AI主导的集中化与板块轮动}

摩根斯坦利：黄金的下一段行情，关键不在央行，而在ETF能否重新入场

当前市场对黄金的讨论，大多集中在央行购金、地缘冲突和通胀对冲这三个传统叙事上。但摩根斯坦利在最新发布的这份报告中，提出了一个更值得推敲的判断：黄金的结构性支撑仍然坚实，但短期内推动价格从当前水平走向每盎司5200美元的关键变量，已经不是央行，而是ETF资金流的重新激活。报告明确指出，如果没有ETF需求的配合，他们的目标价将变得难以实现。

这份报告的价值不在于重申“黄金长期看多”这一共识，而在于它精准地拆解了当前黄金定价中三个相互拉扯的力量：地缘缓和的正面效应、美联储鹰派立场的压制、以及央行结构性买盘的托底。这三个力量叠加的结果，不是简单的多空判断，而是一个需要特定触发条件才能兑现的价格路径。

报告的核心主张可以概括为一句话：黄金的上行风险仍在，但通往5200美元的道路已经变窄，而这条道路能否走通，取决于一个尚未被满足的条件——ETF买家是否愿意重新进场。

很多人直觉上认为，中东冲突缓和会削弱黄金的避险需求，从而打压金价。但摩根斯坦利的分析得出了一个反直觉的结论：冲突缓和反而对黄金构成支撑。原因在于，本轮中东冲突中，黄金并没有扮演传统的避险角色。

报告提供了一个关键的分析框架：供给冲击比需求冲击更难对冲。中东冲突本质上是一个供给侧的冲击——它推高油价、抬高债券收益率、并对油气进口国的财政平衡造成压力。在这种环境下，黄金的避险功能反而受到抑制。报告特别指出，土耳其等国家在压力下甚至不得不卖出黄金。数据显示，冲突升级期间黄金的平均日表现为负值，而缓和期间反而录得正回报。

这意味着，地缘缓和对黄金的正面效应，不是通过“避险需求下降”这个渠道，而是通过“缓解油价-通胀-利率”这个间接链条实现的。摩根斯坦利的石油策略师已经下调了油价预期，这将减轻央行的通胀压力，进而降低它们卖出黄金或收紧货币政策的必要性。从这个角度看，中东缓和是黄金的一个结构性利好，而非利空。

这一层分析对投资者的含义是：不要用传统的“地缘风险=买黄金”的框架来理解当前阶段。黄金与地缘的关系已经发生了结构性变化，其定价权重正在向利率和实际收益率倾斜。

在黄金的所有需求来源中，官方部门的购金行为是最具结构性的力量。世界黄金协会的最新调查显示，45\%的受访央行预计未来12个月将增加黄金储备，这一比例创下历史新高。中国央行已经明显加快了购金步伐——报告数据显示，3月至5月期间中国央行购买了23吨黄金，而此前12个月的总购买量仅为19吨。

但这里需要区分两个概念：央行购金的“存量支撑”和“边际推动力”。央行购金为金价提供了一个坚实的地板，但要让金价从当前水平继续大幅上行，需要的是边际上的增量买家。而央行的购金行为虽然稳定，但其节奏受到外汇储备管理、本币汇率稳定等多重因素制约，很难在短期内出现爆发式增长。

报告也暗示了这一点：尽管央行需求是结构性支撑，但仅凭央行一家之力，难以推动金价突破5200美元。印度的情况提供了另一个视角——印度在5月上调了黄金进口关税以抑制购买，这表明一些新兴市场国家正在主动管理黄金的流入速度。

对于投资者而言，央行购金是一个“持仓不动”的理由，但不是“加仓追高”的依据。真正决定下一阶段金价弹性的，是另一类资金。

3. 美联储是黄金近端最大的约束，而ETF是这一约束的传导载体

摩根斯坦利对美国经济团队的判断非常明确：最新的FOMC声明、点阵图和新闻发布会都偏鹰派，且声明中没有提及劳动力市场的下行风险。这直接推高了市场对加息的预期，并压低了金价。报告中的数据显示，市场定价的加息次数从5月底的0.5次迅速攀升至6月中旬的1.6次。

美联储的鹰派立场对黄金的压制，主要通过一个渠道传导：机会成本。当利率上升时，持有不生息的黄金的机会成本随之增加。但报告进一

[... middle omitted ...]

据显示，黄金在加息后并非必然下跌，但这次的条件不同

报告提供了一个值得注意的历史数据：在25个基点的加息后，黄金一个月平均上涨0.84\%。这个数据本身并不支持“加息必然利空黄金”的简单结论。报告还引用了世界黄金协会的几个历史案例——2006年6月、2017年3月、2018年12月、2022年11月和2023年3月——在这些案例中，加息后黄金都出现了不同程度的上涨。

但这些历史案例的共性是什么？每一次黄金在加息后上涨，都不是因为加息本身，而是因为市场将加息解读为“政策即将见顶”或“政策失误”。2018年12月的加息被市场视为错误，美联储随后在2019年转为降息；2023年3月的加息发生在银行业危机背景下，市场预期政策将很快转向宽松。

而当前的环境与这些历史案例存在一个根本差异：市场并不认为美联储即将转向。摩根斯坦利的经济学家预计，美联储将维持利率不变直至2026年。这意味着，黄金无法从“预期宽松”这个渠道获得支撑。历史案例提供的乐观信号，在当前条件下可能并不适用。

\subsection{[R003] JPM：中东局势缓和，日本日化与造纸股迎来结构性成本重估}
{\small\color{refgray}归类：能源与大宗：油价预期下调，成本传导时滞创造机会 / 日本股市：板块分化加剧，制药出清与消费成本重估 / 风险偏好与资金流向：AI主导的集中化与板块轮动}

地缘政治风险从来不只是新闻标题。当美国与伊朗在6月17日签署谅解备忘录、暂时结束敌对状态的消息传出，全球资本市场的第一反应是油价回落。但JPM日本股票研究团队在最新发布的一份跨行业报告中提出了一个更值得深思的问题：中东局势改善对不同行业、不同公司的利润影响，是否被市场以同一把尺子衡量？答案是否定的。这份报告的核心判断是，市场低估了成本端改善对日本个人护理公司的利润弹性，高估了造纸企业的定价权风险，而对化妆品公司的传导链条则完全看错了方向。

这不是一则简单的“油价跌了，成本降了”的故事。真正需要拆解的是，成本下降的时滞、传导机制、以及不同行业在面对成本改善时截然不同的议价能力。这些变量，决定了哪一类公司能从地缘政治红利中真正获益。

1. 个人护理公司拥有最清晰的成本弹性，但市场可能忽略了4-6个月的传导时滞

JPM分析师Akiko Kuwahara在报告中给出了一个非常具体的测算：原油价格每变动1美元/桶（按年计算），对Unicharm和花王的营业利润影响约为5-6亿日元，对狮王的影响约为1-2亿日元。这个数字本身并不惊人，但关键变量在于传导时滞——报告明确指出，原材料价格波动需要4-6个月才能反映到这些公司的损益表中。

这意味着什么？假设当前油价下跌趋势能够维持，那么对2027财年（日本财年通常从4月开始）的成本压力缓解将是实质性的。但市场往往只关注即期的油价变动，而忽略了利润表反映的滞后性。这恰恰为投资者提供了一个时间窗口：在油价下行趋势确立后，个人护理公司的盈利预期存在上修空间，而这个上修尚未被充分定价。

更重要的是，Unicharm在中东地区有直接业务敞口。局势缓和不仅降低其原材料成本，还消除了其区域销售环境的不确定性。这是一种双重利好：成本端和收入端同时改善。但报告也提醒，2026财年下半年存在一个不可忽视的风险——家庭因恐慌而囤货的需求可能回落，同时前期高价采购的原材料库存成本尚未消化完毕。这意味着，短期利润可能先承压，中期才迎来真正的改善。

2. 造纸企业的成本利好被价格谈判风险对冲，市场对此定价不足

造纸行业的情况要复杂得多。从表面看，油价下跌对纸浆和包装材料的成本压力同样构成利好。JPM测算，原油价格每变动1美元/桶对王子控股和日本制纸的营业利润影响约为3亿日元，大王制纸约1亿日元，连合包装仅约1000万日元。这些弹性数字并不小。

但报告揭示了一个被市场忽视的反向风险：日本主要造纸厂商原计划在7-8月对印刷纸等产品提价。油价突然大幅下跌，反而可能削弱它们在价格谈判中的筹码。客户会问：你的成本都降了，为什么还要涨价？这一逻辑一旦成立，不仅提价计划可能受阻，甚至连已经建立的固定成本转嫁机制都可能被延迟。

这是典型的“成本利好被收入端恶化对冲”的案例。市场如果只看到成本下降就买入造纸股，很可能忽略了价格谈判环境恶化带来的隐性损失。造纸企业的股价反弹，需要等待成本改善真正转化为利润，而不是被价格竞争所吞噬。

化妆品行业是这份报告中最容易被误读的板块。直觉上，中东局势改善对化妆品公司的影响似乎最小——因为包装材料和部分原料的采购成本占比本来就不高，市场此前对这一风险的担忧也相对有限。JPM对此的判断是：直接成本影响确实有限。

但真正的故事不在成本端，而在需求端。报告指出，消费者信心回升和旅游人数增加，才是化妆品公司受益的主要路径。这意味着，这是一条更慢、更间接、但可能更持久的传导链条。化妆品公司不会因为油价跌了一个月就立刻利润大增，但如果中东稳定能够持续提振全球消费者信心、推动国际旅游复苏，那么高端化妆品和旅游零售渠道的销售弹性将逐步显现。

市场目前对化妆品板块的定价，可能仍然停留在“成本影响有限因此无甚影响”的线性思维上，而低估了地缘政治改善对消费场景修复的催化作用。

4. 报

[... middle omitted ...]

反而削弱了提价理由。化妆品公司之所以走另一条路，是因为其成本结构中原料占比低，但需求对宏观环境敏感。

这个框架本身是清晰的，但报告没有给出一个系统的量化模型来评估不同公司的“成本转嫁系数”。例如，Unicharm和花王之间的转嫁能力是否存在差异？狮王的宠物护理业务是否比其个人护理业务更具定价权？这些细分层面的分析，可能需要投资者结合公司自身的品牌组合、渠道结构、客户集中度等变量进一步拆解。

基于以上分析，我们可以提炼出一个三层观察框架，用来判断哪类公司能从地缘政治改善中真正受益：

第一层，看成本弹性。JPM的测算已经提供了基础数据，但投资者需要关注的是弹性是否被重复定价。如果一家公司的成本弹性已经被市场充分计入股价，那么后续的上涨空间就有限。

第二层，看传导时滞。4-6个月的时间差意味着，当前油价下跌对利润的提振最早也要到2026年第四季度才能体现。市场情绪可能在短期内推动股价上涨，但真正的业绩验证需要时间。投资者需要区分“预期驱动”和“业绩驱动”的上涨阶段。

\subsection{[R004] JPM：日本制药板块正在经历的不是回调，而是结构性出清}
{\small\color{refgray}归类：日本股市：板块分化加剧，制药出清与消费成本重估 / 风险偏好与资金流向：AI主导的集中化与板块轮动}

2026年6月，日本股市在AI和半导体板块的带动下创下历史新高，日经225指数在6月19日收于纪录高位。但如果你只看制药板块，会以为自己身处另一个市场。JPM最新一周的日本制药行业周报揭示了一个令人不安但必须正视的事实：日本制药板块正在经历的不是短期回调，而是一场由资金结构、研发风险和估值体系共同驱动的出清过程。

这份报告的核心判断是：日本制药板块的疲软并非周期性波动，而是市场正在用AI时代的估值逻辑重新定价整个行业。那些无法在2026年下半年拿出明确催化剂的公司，将继续被市场抛弃。但报告同时指出，2026年下半年的催化剂事件可能成为板块反转的关键节点。问题在于，这些催化剂是否足够强大，足以逆转当前的结构性压力。

以下是我们从这份报告中提炼的五个关键洞察，以及一个尚未被充分讨论的关键变量。

1. 日本制药板块正在经历的不是行业问题，而是资金流向的结构性挤压

JPM的数据显示，在截至6月19日的一周内，制药板块仅上涨0.8\%，而同期TOPIX指数上涨4.2\%。拉长到一个月维度，制药板块下跌5.6\%，而TOPIX上涨5.0\%。更令人警醒的是，在三个月和六个月的时间跨度内，制药板块分别下跌11.1\%和4.2\%，而TOPIX同期上涨12.1\%和19.5\%。

这不是某个公司的业绩问题，这是整个板块在资金争夺战中的系统性溃败。日本股市的上涨高度集中于AI和半导体相关板块，非铁金属、电子设备、玻璃与陶瓷产品等板块在过去三个月分别上涨28.0\%、44.0\%和38.7\%。相比之下，制药板块在行业分类中排名倒数。

这意味着什么？日本制药公司当前的估值压力，很大程度上是“被动的”——不是因为基本面恶化，而是因为市场可配置资金被AI和半导体主题抽走。对于持有日本制药股的投资者来说，最大的风险可能不是公司管线失败，而是市场根本没有关注管线。

2. 个股分化揭示了一个残酷规律：市场只奖励有明确催化剂的公司

在板块整体疲软的背景下，少数公司跑出了超额收益。住友制药和协和麒麟在当周分别上涨6.9\%和5.7\%，跑赢TOPIX。但这两家公司的上涨逻辑截然不同，且都带有明显的“催化剂驱动”特征。

协和麒麟是过去一个月板块中表现最好的公司，涨幅达6.5\%。但JPM明确给出“UW”（减持）评级，目标价1800日元，较当前2538日元有29.1\%的下行空间。这份报告的逻辑是：尽管rocatinlimab开发终止的坏消息在短期内已被消化，但公司的中长期前景仍然不确定。市场关注的焦点已经转向KHK4951的二期临床数据，该药物针对湿性年龄相关性黄斑变性和糖尿病性黄斑水肿，数据可能在10月的学术会议上公布。

住友制药的情况则更为复杂。JPM给予“OW”（增持）评级，目标价2300日元，隐含56.5\%的上行空间。公司在欧洲血液学协会大会上展示了nuvisertib和enzomenib的最新数据，同时其合作方诺和诺德的Wegovy获得MASH适应症批准。但即使如此，住友制药在三个月和六个月维度分别下跌20.8\%和35.9\%。

这里有一个需要深入思考的问题：为什么一家被分析师看好、有56\%上行空间的公司，在中期维度上仍然被市场大幅抛售？答案可能在于，市场对日本制药公司的估值正在从“管线价值”转向“现金回报价值”。那些需要持续投入研发、短期内无法产生稳定现金流的公司，正在被重新定价。

3. 估值数据揭示了一个被忽视的真相：市场正在用AI标准重新定义“便宜”

JPM提供的估值表格是整份报告中最值得反复研读的部分。截至6月19日，日本主要制药公司的估值水平呈现出极度分化的特征：

中外制药：PE 24.2倍（FY26E），PB 5.4倍，ROE 23.6\%，JPM给予OW评级，目标价11000日元，隐含45.7\%上行空间。 安斯泰来：PE 9.9

[... middle omitted ...]

这个指标。住友制药的ROE高达25.2\%，但其PE仅为6.7倍。这说明市场不相信这个ROE水平能够持续。同样，安斯泰来22.0\%的ROE也只换来了9.9倍的PE。相比之下，中外制药23.6\%的ROE支撑了24.2倍的PE，差距在于市场对中外制药管线可持续性的信心。

对于投资者而言，这意味着在日本制药板块中，“便宜”不再是买入的理由。市场正在用AI行业的估值框架来审视制药公司：只有那些能够证明自己拥有持续创新能力和明确商业化路径的公司，才能获得估值溢价。

4. 2026年下半年催化剂清单是板块反转的必要条件，但并非充分条件

JPM在报告中明确表示，预计制药板块短期内仍将保持疲软，但2026年下半年的催化剂事件可能推动股价上涨。报告提到了几个关键时间节点：

协和麒麟的KHK4951二期数据可能在10月学术会议上公布。 中外制药的Elevidys（DMD基因疗法）说明会已于6月17日举行。 参天制药的Rhopressa（青光眼）已于6月19日获得日本制造和销售批准。

\subsection{[R005] GS：618大促揭示中国美妆市场真正的分水岭，不是国货与外资，而是高端与大众}
{\small\color{refgray}归类：中国消费：618大促揭示高端化结构性力量 / 风险偏好与资金流向：AI主导的集中化与板块轮动}

GS：618大促揭示中国美妆市场真正的分水岭，不是国货与外资，而是高端与大众

市场对2026年618美妆大促的普遍解读，往往落入“国货崛起”或“外资反攻”的二元叙事。但GS这份最新发布的618抖音渠道脉冲检查报告，揭示了一个更深层、也更值得产业决策者关注的信号：真正的分水岭不在于品牌国籍，而在于价格带。高端化正在从趋势变为结构性力量，而这一力量的受益者，并非所有高端品牌。

这份报告的数据窗口覆盖2026年5月1日至6月18日，聚焦抖音平台。抖音作为当前中国美妆线上增长最快的渠道，其表现往往领先于天猫和京东，成为品牌战略调整的先行指标。GS的追踪数据显示，抖音美妆Top 500品牌整体GMV同比增长约35\%，但这一数字背后，是极度分化的格局：头部50个品牌增速高达约40\%，而排名100名之后的品牌增速已回落至20\%左右。市场正在加速集中，而集中的方向，是高端。

1. 外资高端品牌在抖音的“反攻”并非全面胜利，而是集团内部资源向高利润线集中

GS数据显示，在抖音Top 500品牌中，外资品牌在各个价格带的完成率普遍高于国产品牌，平均完成率约140\%，而国产品牌约110\%。但“外资”这个标签过于笼统。真正驱动外资整体表现的，是雅诗兰黛集团和欧莱雅集团旗下高端线的强势。

雅诗兰黛集团在抖音的2026年5月1日至6月18日GMV同比增长56\%，其中雅诗兰黛主品牌增长65\%，海蓝之谜增长54\%，两者完成率分别高达167\%和154\%。欧莱雅集团整体增长33\%，但驱动来自修丽可（增长89\%，完成率190\%）和YSL（增长74\%，完成率176\%），而大众品牌巴黎欧莱雅仅增长7\%，美宝莲甚至下滑7\%。

这意味着什么？外资巨头在抖音的增长，并非“全线复苏”，而是集团战略性地将资源、流量和折扣向高利润、高客单价的品牌倾斜。大众线在抖音的衰退，是集团主动选择的结果，而非被动失利。对于投资者而言，判断一家外资美妆集团在中国的前景，不应看其整体增速，而应拆解其高端线占比的变化趋势。高端线占比持续提升的公司，其盈利能力和品牌护城河才是真正在加强的。

2. 毛戈平是本次618最大的结构性赢家，但它的成功逻辑并非简单的“国货替代”

在所有覆盖的国产品牌中，毛戈平（MGP）的表现最为亮眼：抖音全周期GMV同比增长54\%，6月1日至18日增速更是高达91\%，完成率达到150\%。GS在报告中明确将毛戈平列为首选买入标的，并指出其5月销售偏弱是公司主动的节奏控制，与双十一等大促期间的表现模式一致。

毛戈平的成功，不能简单地归因于“国货替代”或“情绪消费”。它的定价处于中高端价格带，核心单品售价在260-400元区间，直接与外资中高端品牌竞争。在抖音这个以低价和冲动消费闻名的渠道，毛戈平能够实现90\%以上的6月增速，说明它已经建立起超越渠道特性的品牌认知和用户粘性。

更值得关注的是，在GS的分类中，国产高端品牌（如毛戈平）的线上增速，在2026年已经持续跑赢外资高端品牌。Exhibit 3的数据显示，国产高端品牌在2026年3月至5月的线上GMV增速分别为110\%、40\%和40\%，而外资高端品牌同期仅为20\%、20\%和30\%。这意味着，在高端化这条赛道上，中国品牌终于出现了具备结构性竞争力的玩家，而不再是仅仅依靠性价比在低端市场厮杀。

3. 大众市场正在经历“无声的失血”，国产品牌的规模优势正在被侵蚀

与高端市场的繁荣形成鲜明对比的，是大众市场的困境。GS的数据显示，在抖音Top 500品牌中，排名100名之后的品牌增速已降至20\%左右，显著低于头部品牌。而具体到国产品牌，珀莱雅在5月1日至6月18日期间抖音GMV同比增长24\%，但6月1日至18日已经转为下滑5\%，完成率121\%。上海家化虽然整体增长68\%，但6月1日至18日下滑1

[... middle omitted ...]

力度均小于去年），依赖低价和流量投喂的增长模式便难以为继。对于产业决策者而言，这意味着过去几年“所有品牌都能在抖音增长”的窗口期已经关闭。大众品牌如果不能快速向中高端迁移，或者找到极致的成本效率优势，将面临持续的份额流失。

4. 报告尚未完全回答的关键问题：天猫渠道的表现和补贴的真实影响

GS这份报告聚焦于抖音渠道，但明确指出，618大促中天猫和京东渠道同样重要。报告提到，618期间（5月+6月）的GMV约占二季度线上GMV的75\%，占全年GMV的21\%。抖音只是其中一个战场，天猫仍然是美妆品牌最大的线上阵地。

报告没有完全展开的几个关键问题包括：第一，天猫渠道的表现是否与抖音一致？尤其是毛戈平在抖音的高增长，是否以牺牲天猫渠道为代价？第二，报告提到“平台和政府的补贴”可能影响品牌的实际收入和利润率，但抖音数据无法反映折扣深度和退货率。第三，抖音的“完成率”指标是基于2025年同期数据，而2025年618的基数效应可能因疫情等因素存在扭曲，这一点需要进一步验证。

\section{Disclaimer}
{\scriptsize\color{refgray}
This community is a paid learning and information-sharing community created for general educational purposes only. The membership fee is charged solely for access to community discussions, educational content, organization, and operational costs. It is not an advisory fee, management fee, performance fee, brokerage fee, or compensation for personalized financial, investment, legal, tax, or professional advice.\par
I participate and share information solely as an individual and for educational discussion among community members. I am not acting as a financial adviser, investment adviser, broker, dealer, portfolio manager, legal adviser, tax adviser, accountant, fiduciary, or any other licensed professional adviser. Nothing shared in this community constitutes, and should not be interpreted as, financial advice, investment advice, legal advice, tax advice, a recommendation, solicitation, offer, endorsement, instruction, or guarantee to buy, sell, hold, short, trade, or invest in any security, cryptocurrency, fund, derivative, strategy, or other financial product.\par
All content shared in this community, including messages, discussions, links, files, charts, examples, market commentary, watchlists, AI-generated or AI-polished summaries, and third-party materials, is general, impersonal, and educational in nature. It does not take into account any individual's financial situation, investment objectives, risk tolerance, investment horizon, liquidity needs, tax status, legal restrictions, or suitability requirements. No content should be relied upon as suitable for any specific person, account, portfolio, or financial decision.\par
Information shared in this community may be incomplete, inaccurate, outdated, speculative, AI-generated, AI-edited, or based on third-party internet sources that have not been independently verified. I make no representation or warranty as to the accuracy, completeness, timeliness, reliability, or suitability of any information shared.\par
Financial markets involve substantial risk, including the possible loss of principal. Past performance, historical data, hypothetical examples, simulated results, backtests, personal opinions, or educational examples do not guarantee future results. Each member is solely responsible for conducting their own research, exercising independent judgment, and making their own decisions. Before making any financial, investment, legal, or tax decision, members should consult qualified and licensed professionals.\par
Participation in this community, payment of any membership fee, or communication with me or other members does not create any adviser-client, fiduciary, professional, contractual, agency, partnership, employment, or other advisory relationship. By joining or participating in this community, you acknowledge and agree that you are solely responsible for your own decisions, actions, transactions, profits, losses, risks, and consequences, and that I shall not be liable for any direct, indirect, incidental, consequential, financial, legal, tax, or other loss or damage arising from or related to any content, discussion, information, or material shared in this community.\par
}
\end{document}
